% !TeX root = ../se200_referential_regimes.tex

% =============================================================================
\section{Conclusion}
\label{sec:conclusion}
% =============================================================================

Identity in a neutral substrate is not fixed by naming kinds alone.
It is fixed by specifying which transformations preserve sameness.
Once the ordinary operations of accountability systems are specified,
individuation becomes a structural question rather than a taxonomic preference.

This paper began with the neutral-substrate constraint developed in
\emph{Neutral Substrates}~\citep{case2026neutral}.
A substrate that must remain usable under persistent disagreement
may commit to stable reference and permitted attribution,
but not to object-level causal or normative interpretation.
This paper explores what identity structure that constraint requires
when records are transformed in ordinary use.

The answer is a lower-bound core.
An accountability substrate must be able to refer to six kinds of things:
obligation-bearers, rules, occurrences, scopes, records, and plain referents.
Those six reference kinds supply carriers, but do not settle identity.
Three of them are assigned more than one admissible basis:
a scope by extension or by structure,
a rule by content or by structure,
a plain referent by locus or by object.
Ordinary transformations force apart those alternatives:
decomposition separates extension-fixed from structure-fixed scopes,
refinement separates content-fixed from structure-fixed rules,
and forking separates locus-fixed from object-fixed plain referents.
Thus six reference kinds force at least nine identity regimes:
\[
\begin{aligned}
&\OBL,\quad \OCC,\quad \REC,\quad \LOC,\quad \OBJ,\\
&\SCOPEE,\quad \SCOPES,\quad \RULEC,\quad \RULES .
\end{aligned}
\]

The central result is the identification of the additional regimes.
They are not additional reference kinds, but distinct identity bases forced apart by transformation behavior.
An analysis that names kinds may undercount the structure required for stable reference.
Collapsing regimes does not produce genuine neutrality.
A collapsed distinction must change identity behavior
or be recovered through a hidden discriminator.

The result is conditional.
It holds under the assumptions fixed in this paper:
persistent disagreement, strong stable reference, minimality,
and exclusion of hidden regimes.
It is a lower bound rather than an upper bound or a completeness claim.
A richer transformation basis may leave it unchanged or force additional regimes.

The result is diagnostic rather than prescriptive.
It does not claim that every ontology should expose these nine regimes,
nor that contextual typing, role predicates, or hidden discriminators
are useless as engineering devices.
It says that such devices cannot substitute for regime-level identity commitments
in a neutral substrate under persistent disagreement.
When stable reference must be held as common ground before
interpretation is settled, the substrate structure must preserve
the relevant identity distinctions.

Three papers in this series explore the implications of that constraint.
\emph{Neutral Substrates} states the neutrality constraint.
\emph{Referential Regimes} measures its identity cost.
\emph{Accountable Records} studies how that cost can be made checkable in deployed record systems.
This work addresses the middle step:
stable reference under persistent disagreement requires
transformation-invariant identity regimes,
and at least nine such regimes are forced by the ordinary operations
accountability systems already perform.
